\section{Scanning as modular flow: from rotation algebras to modular dynamics}
\label{sec:scanning_as_modular_flow}

This section explains how the scan algebra (O6) can be viewed as a minimal cross-section of modular dynamics. The main point is structural: the Weyl relation singles out the irrational rotation algebra (a noncommutative torus), while the modular generators $S,T$ provide canonical geometric symmetries whose cusp action supplies a scale-inversion template.

\subsection{The rotation algebra \texorpdfstring{$A_\alpha$}{A\_alpha} as the minimal scan closure}
\label{sec:rotation_algebra}

The $C^*$-algebra generated by unitaries $U,V$ satisfying
\begin{equation}
UV=\e^{2\pi\iu\alpha}VU,\qquad \alpha\notin\QQ,
\end{equation}
is the (irrational) rotation algebra $A_\alpha$ \cite{Rieffel1981,Connes1994}. In the scan interpretation:
\begin{itemize}
  \item $U$ encodes the tick iteration (time as an integer counter),
  \item $V$ encodes the phase/readout coordinate.
\end{itemize}
This is the smallest closed algebraic setting in which scanning and phase readout coexist while remaining intrinsically noncommutative.

\subsection{\texorpdfstring{$T$}{T}-translation and phase periodicity}
\label{sec:T_translation}

On $\HH$, the modular translation $T:\tau\mapsto\tau+1$ induces
\begin{equation}
q = \e^{2\pi \iu \tau} \mapsto \e^{2\pi \iu (\tau+1)} = q,
\end{equation}
so the $q$-coordinate is $T$-periodic. At the same time, the real part shifts by an integer:
\begin{equation}
x := \Re\tau \pmod 1 \quad \Rightarrow \quad x \mapsto x+1 \equiv x.
\end{equation}
This identifies a canonical boundary phase coordinate $x\in\RR/\ZZ$ compatible with modular periodicity. The scan orbit of O5,
\begin{equation}
x_t = x_0 + t\alpha \pmod 1,
\end{equation}
is then an irrational translation on the same circle: it respects the same quotient structure but selects an irrational slope (a minimal, uniquely ergodic orbit).

\subsection{\texorpdfstring{$S$}{S}-inversion, cusp equivalence, and a scale-exchange template}
\label{sec:S_inversion}

The modular inversion $S:\tau\mapsto -1/\tau$ exchanges deep and shallow regions of the cusp. Indeed, if $\Im\tau$ is large (near $\iu\infty$), then $\Im(-1/\tau)$ is small (near $0$). Since $0\sim\infty$ as cusps on $X(1)$ (Section~\ref{sec:modular_group}), $S$ supplies a canonical endpoint identification on the arithmetic quotient.

\begin{remark}[interpretation-layer template]
One may use $S$ as a geometric prototype for ``scale exchange'' (UV/IR) when translating protocol structures into physical semantics. This is a Layer~2 reading and is not required for the Layer~0/1 closure.
\end{remark}


